\documentclass{article}
\title{ベイズ統計入門 — 確率を更新する考え方}
\author{TeX 図書館}

\begin{document}

\section{この教科書のために — 頻度論とベイズ}

「コインが表になる確率は 0.5」と言ったとき、その 0.5 は何を意味しているでしょうか。

\begin{itemize}
\item \textbf{頻度論}の立場: 確率は「無限回繰り返したときの頻度」。パラメータ(真の確率)は\textbf{固定された定数}で、データがその推定値に誤差を与える。検定・信頼区間はこの立場。
\item \textbf{ベイズ}の立場: 確率は「信念の度合い」。パラメータ自体が\textbf{確率分布}をもち、データを得るたびにその分布を\textbf{更新}する。
\end{itemize}

この教科書はベイズの立場を学びます。核となる仕掛けはただ一つ、\textbf{ベイズの定理}です(『統計学入門』第 4 節で出会った式の再登場です)。

\begin{quote}
\textbf{【心構え】}\ ベイズは「結果が出たら考えを更新する」ことの数学化であり、要素技術は中学レベルの条件付き確率だけです。積分の概念(分布の面積を 1 に正規化する)を使いますが、式変形は最小限に抑えます。
\end{quote}

\section{条件付き確率からベイズの定理へ}

復習です。事象 \(B\) が起きたという条件のもとで \(A\) が起きる確率は
\[ P(A \mid B) = \frac{P(A \cap B)}{P(B)} \]
でした。ここで \(A\) と \(B\) の役割を入れ替えて \(P(B \mid A)\) と組み合わせると
\[ P(A \mid B) = \frac{P(B \mid A)\, P(A)}{P(B)} \]
が得られます。これが\textbf{ベイズの定理}です。

\begin{quote}
\textbf{【例:病気の検査】}\ 有病率 0.1\%、感度 99\%、偽陽性率 5\% の検査で陽性が出たとき、実際に病気である確率は
\[ P(\text{病} \mid \text{陽}) = \frac{0.99 \times 0.001}{0.99 \times 0.001 + 0.05 \times 0.999} \approx 1.9\% \]
でした。まれな事象では偽陽性が圧倒する——ベイズ統計は、この「条件の向きを正しく直す」操作を\textbf{一般のパラメータ推定}に拡張したものです。
\end{quote}

\begin{quote}
\textbf{【演習】}\ 上の検査で、事前に「家族に病歴があるので有病率は 10\% と考えられる」ケースの事後確率を計算せよ。
\end{quote}

\textbf{解答の指針:}\ \(\frac{0.99 \times 0.1}{0.99 \times 0.1 + 0.05 \times 0.9} \approx 68.8\%\)。\textbf{事前の信念が結果を大きく変える}——これがベイズの本質です。

\section{事前分布・尤度・事後分布}

連続的なパラメータに一般化します。推定したい量(たとえば「コインの表が出る真の確率」\(\theta\))について:

\begin{enumerate}
\item \textbf{事前分布} \(p(\theta)\): データを見る前の信念。\(\theta\) がどんな値にありそうかの分布。
\item \textbf{尤度} \(p(D \mid \theta)\): 「\(\theta\) が \(x\) だったとしたら、このデータ \(D\) が観測される確率」。データ側の声。
\item \textbf{事後分布} \(p(\theta \mid D)\): データを見たあとの更新された信念。
\end{enumerate}

そして関係式は\textbf{比例形}で十分です:
\[ p(\theta \mid D) \;\propto\; p(D \mid \theta) \, p(\theta) \]
「\textbf{事後 ∝ 尤度 × 事前}」。分母(周辺確率)は \(\theta\) によらず定数なので、形を知るだけなら不要です。

\begin{quote}
\textbf{【例】}\ 「工場の不良品率 \(\theta\) はたぶん 1\% 前後だ」という事前分布があり、100 個検査して 3 個不良だった。尤度は「二項分布で 100 回中 3 回不良」の確率。事後分布は「1\% 前後だった信念が、3 個というデータでやや右に引っ張られた形」になります。
\end{quote}

\begin{quote}
\textbf{【演習】}\ 「事後 ∝ 尤度 × 事前」において、データ件数が増えるとどちらの項が支配的になるか考察せよ。
\end{quote}

\textbf{解答の指針:}\ データが増えるほど尤度は鋭くなり、事前の形の影響は相対的に薄れる。\textbf{データが十分なら事前の選び方は重要でなくなる}——ベイズの重要な性質です。

\section{ベイズ更新の連鎖 — 今日の事後は明日の事前}

ベイズの美点は\textbf{更新を何度も繰り返せる}ことです。今日計算した事後分布は、明日データが追加されたときの事前分布として使えます。

\begin{quote}
\textbf{【例:コインの連続更新】}\ 「表が出る確率 \(\theta\) は不明」という完全な無知から始めます。
\begin{enumerate}
\item 1 回投げて表 → \(\theta\) はやや大きめかもしれない、という事後分布
\item その分布を事前にして、2 回目: 表 → さらに右へ
\item 3 回目: 裏 → 少し左へ
\item …と、データのたびに分布が動いていく
\end{enumerate}
初期の動きは激しく、データが溜まるほど動きは小さくなります(収束)。
\end{quote}

頻度論の「点推定が 1 回で確定する」のに対し、ベイズは\textbf{信念の軌跡}を保持します。意思決定(たとえば工場の検査体制を変えるか)を段階的に行う現場では、この軌跡そのものが役立ちます。

\begin{quote}
\textbf{【演習】}\ A/B テストで「新しいボタンのクリック率」をベイズ更新する設計を、事前・データ・事後の 3 つで記述せよ。
\end{quote}

\textbf{解答の指針:}\ 事前 = クリック率の信念(旧ボタンの実績から)。データ = 新ボタンが表示された回数とクリック回数。事後 = 更新された信念。表示件数が増えるごとに更新を続け、「新ボタンが旧ボタンを上回る確率」を随時計算できる。

\section{共役事前分布 — 手計算できる幸運な組}

事後分布が解析的に計算できる「相性の良い」事前分布の組を\textbf{共役事前分布}といいます。代表的なのが\textbf{ベータ分布}と二項尤度の組です。

コイン(表の確率 \(\theta\))について、事前分布としてベータ分布
\[ \mathrm{Beta}(\alpha, \beta) \]
を使うと、\(n\) 回中 \(k\) 回表というデータに対する事後分布は
\[ \mathrm{Beta}(\alpha + k,\ \beta + n - k) \]
——\textbf{足し算だけで更新が完了}します。\(\alpha\) は「仮想的な表の回数」、\(\beta\) は「仮想的な裏の回数」と読むと直感的です。

\begin{quote}
\textbf{【例】}\ 事前 \(\mathrm{Beta}(2, 2)\)(「たぶん 0.5 前後、弱い信念」)に、3 回中 2 回表のデータ。\(\mathrm{Beta}(4, 3)\) が事後で、事後平均は \(4/7 \approx 0.57\)。事前平均 0.5 がデータ(表優位)で右へ引かれました。
\end{quote}

事前の\textbf{強さ}も重要です。\(\mathrm{Beta}(100, 100)\) は「ほぼ 0.5 に違いない」という\textbf{強い}信念(仮想データ 200 回ぶん)で、同じデータでも事後はほとんど動きません。一方 \(\mathrm{Beta}(1, 1)\)(一様分布)は完全な無知で、データに忠実です。

\begin{quote}
\textbf{【演習】}\ 事前 \(\mathrm{Beta}(50, 50)\) に「10 回中 9 回表」のデータが来た。事後分布と事後平均を求めよ。
\end{quote}

\textbf{解答の指針:}\ \(\mathrm{Beta}(59, 51)\)、事後平均 \(59/110 \approx 0.536\)。\textbf{9/10 = 0.9 というデータに比べて事後は 0.5 の側に強く引かれる}——強い事前の重みです。

\section{事後分布からの読み取り — 点推定と信用区間}

事後分布が手に入ったら、次のように読み取ります。

\begin{itemize}
\item \textbf{事後平均・事後中央値}: パラメータの点推定。
\item \textbf{信用区間}(credible interval): 事後分布の中央 95\% の範囲。「\(\theta\) がこの区間に入っている\textbf{確率が 95\%}」と\textbf{直感的に読める}(頻度論の信頼区間は「区間を作り続けたら 95\% 含む」で意味がねじれていたのと対照的)。
\item \textbf{事後確率そのもの}: 「\(A\) ボタンのクリック率が \(B\) を上回る確率は 87\%」のような、意思決定に直接使える数字。
\end{itemize}

\begin{quote}
\textbf{【例】}\ \(\mathrm{Beta}(4, 3)\) の事後で「\(\theta > 0.5\) の確率」を数値積分すると約 0.57。つまり「このコインは表が有利だと考えるべき確率は 57\%」——まだ半信半疑という数字がそのまま出てきます。
\end{quote}

\begin{quote}
\textbf{【演習】}\ 「新薬の治癒率が旧薬より高い事後確率が 96\%」というベイズ解析の結果と、「p 値 0.04」の頻度論の結果を、意思決定者への説明の観点から比較せよ。
\end{quote}

\textbf{解答の指針:}\ ベイズは「高いと考えられる確率 96\%」と直截に読める。p 値は「差がないと仮定したときの極端さ」で、非専門家には誤解されやすい。ただしベイズは事前の選択が結果に影響する点を必ず併記する。

\section{MCMC の直感 — 計算機で事後分布を歩く}

共役事前分布の幸運な組から外れると、事後分布は\textbf{数式で書けません}。そこで計算機の出番です。\textbf{MCMC}(マルコフ連鎖モンテカルロ)は、事後分布を「歩き回って標本を集める」手法の総称です。

発想は次のとおりです。

\begin{enumerate}
\item パラメータ空間のどこかに現在位置がある
\item 近くの候補地点に移動を提案し、「そこでの尤度 × 事前」が現在より良ければ移動、悪ければ\textbf{確率的に}移動(メトロポリス・ヘイスティングスの核)
\item これを何万回と繰り返すと、足取りは「事後分布に比例した密度」で集まる
\item 集まった標本で平均・区間・確率を全部計算できる
\end{enumerate}

\begin{lstlisting}[language=Python]
# メトロポリス法で事後分布からサンプリングする雰囲気
import random, math

def log_posterior(theta, k, n):
    if not 0 < theta < 1:
        return -1e9
    # 二項尤度 × 一様事前(対数)
    return k * math.log(theta) + (n - k) * math.log(1 - theta)

samples, theta, k, n = [], 0.5, 2, 3   # 3回中2回表
for _ in range(50000):
    prop = theta + random.gauss(0, 0.1)      # 近くに移動を提案
    if math.log(random.random()) < log_posterior(prop, k, n) - log_posterior(theta, k, n):
        theta = prop                          # 受理
    samples.append(theta)

mean = sum(samples[5000:]) / len(samples[5000:])   # ウォームアップを除く
print(f"事後平均 ≈ {mean:.3f}")     # 約 0.57
\end{lstlisting}

\begin{quote}
\textbf{【演習】}\ 上のコードで提案幅 0.1 を 0.001 にしたら、サンプリングに何が起きるか予想せよ。
\end{quote}

\textbf{解答の指針:}\ 動きが細かすぎて、事後分布の一部しか探索できない(収束が極端に遅い)。提案幅は「探索の効率」を決める調整項です。

\section{実践 — 不良品率の推定を通しで行う}

小さな工場の例を通しで考えます。過去の実績から「不良品率は 2\% 前後」という事前(\(\mathrm{Beta}(2, 98)\): 仮想 100 個中 2 個不良)を置きます。今日、200 個検査して 7 個不良でした。

\begin{enumerate}
\item \textbf{事後分布}: \(\mathrm{Beta}(2 + 7,\ 98 + 193) = \mathrm{Beta}(9, 291)\)
\item \textbf{事後平均}: \(9/300 = 3.0\%\)——事前の 2\% がデータ(7/200 = 3.5\%)で押し上げられた
\item \textbf{95\% 信用区間}: おおよそ \(1.4\%\)〜\(5.2\%\)。不良品率が 5\% の工程基準を超える確率は約 3\%
\item \textbf{意思決定}: 「基準超過の確率 3\% なら見守る」「8\% が許容上限なら即調査」——事後確率は判断基準に直接掛けられる
\end{enumerate}

\begin{quote}
\textbf{【演習】}\ 上の例で「明日も 200 個検査する」ときの事前分布は何か。また、明日 2 個不良だったら事後平均はどうなるか計算せよ。
\end{quote}

\textbf{解答の指針:}\ 明日の事前は今日の事後 \(\mathrm{Beta}(9, 291)\)。2 個不良なら \(\mathrm{Beta}(11, 489)\)、事後平均 \(11/500 = 2.2\%\)。\textbf{連続更新で信念が安定していく}様子が読み取れる。

\section{おわりに — 信念を更新する作法}

この教科書の骨格を振り返ります。

\textbf{事後 ∝ 尤度 × 事前}という一つの式に、すべてが詰まっています。事前の選び方(無情報にするか実績にするか)、尤度の設計(どの確率モデルがデータの生成過程に合うか)、事後の読み方(平均・信用区間・確率)、そして計算(MCMC)。頻度論と対比すると、ベイズは「\textbf{意思決定者の信念}」を正面から扱う統計です。

次の一歩は、 Stan や PyMC といった本格的な MCMC ツールでのモデリング、そして『確率過程入門』(本図書館)へ進むと、時間とともに変化するパラメータのベイズ推定(状態空間モデル)に繋がります。

\begin{quote}
\textbf{【総合演習】}
\begin{enumerate}
\item 事前 \(\mathrm{Beta}(1, 1)\) に「5 回中 0 回表」のデータ。事後分布と事後平均を求めよ。
\item 「すべての工学部生は数学が得意」という強い事前信念があるとき、反例を 1 つ見ただけで信念は崩れるか。ベイズ更新の観点から答えよ。
\item ベイズ信用区間と頻度論の信頼区間の解釈の違いを、1 文ずつで述べよ。
\end{enumerate}
\end{quote}

\textbf{解答の指針:}\ 1. \(\mathrm{Beta}(1, 5)\)、事後平均 \(1/6 \approx 0.17\)。「0 回表」というデータでも θ=0 とは断定しないのがベイズ(事後分布は 0 側に寄るが幅をもつ)。2. 反例は尤度を大きく動かすが、事前の強さ(たとえば Beta(1000,100))なら事後はほとんど動かない。信念を強く持ちすぎるとデータに耳を借りなくなる。3. 信用区間 = 「θ が区間にある確率 95\%」/ 信頼区間 = 「同じ手順で作った区間の 95\% が真値を含む」。

\end{document}
